\documentclass[11pt]{article}
\usepackage[a4paper,margin=1in]{geometry}
\usepackage{amsmath,amssymb,amsthm}
\usepackage{booktabs}
\usepackage{graphicx}
\usepackage{xcolor}
\usepackage{hyperref}
\hypersetup{colorlinks=true,linkcolor=blue,urlcolor=blue,citecolor=blue}
\newcommand{\rt}[1]{\texttt{#1}}
\newcommand{\fn}[1]{\texttt{#1}}
\newcommand{\oeis}[1]{\href{https://oeis.org/#1}{\texttt{#1}}}
\newcommand{\iftable}[1]{\IfFileExists{#1}{\input{#1}}%
  {\fbox{\parbox{0.9\textwidth}{\centering\small table
  \texttt{\detokenize{#1}} pending}}}}

\theoremstyle{plain}
\newtheorem{proposition}{Proposition}

\title{\textbf{Report R20 --- the growth sequences have names:
twenty-three OEIS identifications}\\
\large Krylov sector analysis of quantum cellular automata}
\author{Tier 1 analysis}
\date{2026-08-14}

\begin{document}
\maketitle

\begin{abstract}
Every growth law in this project was derived here --- from an exact linear
recurrence over the wall grammar (R18), from a closed form (R19), or, in the
cases R16 flagged, from a fit. That fixes an exponent but says nothing about
whether the underlying integer sequence is \emph{known}. We asked the On-Line
Encyclopedia of Integer Sequences: \textbf{24 identifications, covering 38
(rule, observable) pairs and drawing on 18 distinct entries}, each matched as an
exact contiguous block of $17$ terms with the index shift recorded. Three of the
identifications change what we can say. (i) The sequence that \emph{falsifies}
the task's $\varphi$ prediction for rule 28 is \oeis{A023434}, ``dying
rabbits'' --- Fibonacci's rabbits with a finite lifetime --- so the falsification
is a classical combinatorial object rather than a numerical accident, and rule
70 lands on the same entry by independent measurement. (ii) $W_{156}/W_{198}$'s
$D_{\max}$ is \oeis{A178715}, the ``Select All, Copy, Paste'' problem, whose cost
structure (a factor $m$ costs $m+1$) is exactly why the base is $4^{1/5}$, and
whose own recurrence $a(n)=4a(n-5)$ reproduces our ladder term for term.
(iii) $W_{108}$'s frozen-state count is \oeis{A006498}, whose entry states the
Fibonacci-product factorisation R19 derived. Two further identifications are
new results in their own right: $W_{105}$'s $D_{\max}$ is the central binomial
coefficient \oeis{A001405}, and $W_{150}$'s is \oeis{A214282}, the largest Euler
characteristic of a downset on the $N$-cube. We also record what is
\emph{absent}: the four \rt{pbc} series of the dissipative $\rho$-group return
no hits at all, so they appear to be new.
\end{abstract}

\tableofcontents

\section{Why bother}

An OEIS identification is not a theorem, and none of the physics below depends
on one. It buys three specific things.

\paragraph{An external oracle.} An entry carries a b-file with hundreds of
terms. Our frontier stops at $N=17$ for the dissipative rules and $N=22$ for the
unitary ones; the named sequence does not. When the planned extension runs on
larger hardware, the first hundred terms of the continuation are already
written down by someone who has never seen this code. That is the only kind of
check on a growth law that is not produced by the thing it checks.

\paragraph{A name for a falsification.} R9's task sec.~6 predicted
$a_{\rm wcc}\ge\varphi$ for rule 28, on the grounds that it inherits rule 156's
\rt{01} wall grammar. It does not hold, and \fn{test\_wcc.py}'s\linebreak[1]
\fn{test\_rule\_28\_sector\_count\_}\-\fn{follows\_the\_plastic\_number\_not\_phi}
pins why: the \rt{obc0} sector count obeys the exact recurrence
\[
a(N) = a(N-1) + a(N-2) - a(N-4),
\]
whose characteristic polynomial factors as $(x-1)(x^3-x-1)$, giving the plastic
number $\rho = 1.324718$ --- well below $\varphi=1.618034$. That recurrence is
the \emph{definition} of \oeis{A023434}. Section~\ref{sec:rabbits} takes this
further than a coincidence.

\paragraph{A mechanism for an exponent.} $4^{1/5}=1.319508$ is an odd-looking
constant to fall out of a cellular automaton. \oeis{A178715} explains it, and
Section~\ref{sec:copypaste} spells the explanation out.

\section{Method, and what ``matched'' means}

\subsection{Our side}

Unitary series come from \fn{scaling.paper\_figures}'s merged Tier-1a $+$
Tier-1e loader, $N=6\ldots22$ (R19). Dissipative series are assembled by
\fn{scaling.oeis.diss\_series}: the Tier-1a archive for $N\ge6$, and a fresh
recomputation from \fn{graph.scc} for $N\le12$, with the overlap
\emph{asserted} to agree rather than assumed. The recomputation is what supplies
$N=1\ldots5$, and those low terms are exactly what fixes an index shift, so they
had to come from a code path whose agreement with the archive is checked.

Two conventions matter. The ergodic early-exit is switched off during the
recomputation, so that a small unit reports its class count instead of bailing
out --- with it on, $W_{73}$ reads $0,0,0,4,\ldots$ instead of $1,2,3,4,\ldots$
and the shift comes out wrong by three. And $D_{\max}$ for a dissipative rule
means the largest \emph{recurrent} class, i.e.\ the largest key of
\rt{size\_hist}, not the largest weakly connected component; the two differ (for
$W_{70}$ at $N=12$ they are $16$ and $486$) and only the former is the object
R-T14 says the renewal counts.

\subsection{No sequence is inherited from a symmetry partner}

Several identifications cover more than one rule. In every such case the rules
were measured separately and found to agree; none is copied across a symmetry.
This is not pedantry. \fn{core.rules} warns that reflection is \emph{not} a
guaranteed sector-size symmetry for dissipative rules under the even-first
convention, so $W_{28}=W_{70}$ and $W_{29}=W_{71}=W_{157}=W_{199}$ are
measurements, not shortcuts.
\fn{test\_oeis.py::test\_grouped\_rules\_were\_each\_measured\_not\_inherited}
recomputes each member of each group and compares.

\subsection{The OEIS side}

Entries are fetched once and cached in \fn{analytics/oeis\_terms.json}, so the
tests run offline and a failure means our series moved rather than that OEIS
did. A match requires that \emph{every} term we have appears \emph{consecutively
and in order} inside the entry's listed terms; nothing is matched on a truncated
prefix, up to a constant factor, or on a subsequence. The shift is then read off
against the entry's declared offset and stored, so that
\[
\text{ours}(N) \;=\; \texttt{Axxxxxx}(N + \text{shift})
\]
is a checkable statement and not a gesture. All 23 currently verify on 17 terms
each.

\section{The identifications}

\begin{center}
\small
\iftable{tab_r20_oeis.tex}
\end{center}

\noindent and the entries themselves, with the closed form in our own index
where we have one:

\begin{center}
\small
\iftable{tab_r20_oeis_names.tex}
\end{center}

\section{Three that are more than bookkeeping}

\subsection{Dying rabbits, and rule 28}
\label{sec:rabbits}

Rule 28 is \rt{IIVD}. Its \rt{obc0} recurrent-class count and its $D_{\max}$
ladder are \emph{term for term identical to rule 70's} over the whole computed
range, $N=1\ldots17$, at both boundary conditions:
\[
2,3,4,6,8,11,15,20,27,36,48,64,85,113,150,199,264
\]
\[
D_{\max} = 1,2,2,2,4,4,4,8,8,8,16,16,16,32,32,32,64 \;=\; 2^{\lfloor (N+1)/3\rfloor}.
\]
The two rules are reflection partners, so the agreement is unsurprising; it is
nonetheless \emph{measured}, for the reason given above.

The class count is \oeis{A023434}$(N+2)$, ``dying rabbits'':
$a(n)=a(n-1)+a(n-2)-a(n-4)$ --- Fibonacci's rabbits with a finite lifetime, the
$-a(n-4)$ being the deaths. For $W_{28}$ this is not merely the right
recurrence; it reads as the right \emph{story}. Rule 28 inherits rule 156's
\rt{01} wall grammar, which on its own gives Fibonacci and $\varphi$ (and does,
for $W_{156}$: see the table). What rule 28 adds is a \rt{D} channel, a reset to
$|0\rangle$, which destroys a wall after finitely many steps. A wall with a
finite lifetime is precisely what turns Fibonacci's recursion into the dying-
rabbit one, and $\varphi=1.618$ into $\rho=1.325$.

We flag this as an \emph{interpretation to be checked against the wall grammar
in Tier-2}, not a result. What is a result is the arithmetic: the exact
recurrence, the identification, and the fact that the sequence which falsifies
the $\varphi$ prediction is a named classical object with a b-file.

A second observation belongs here. For $W_{28}$ the weakly connected component
count and the recurrent-class count coincide ($11,15,20,27,\ldots$ for both), so
one identification covers both observables. That is not true across the whole
$\rho$-group: $W_{29}$ and $W_{71}$ are the two rules whose WCC count differs
from their recurrent-class count, which is why \fn{sectors.rule\_point} records
$W_{70}$'s base as an ``integer recurrence'' at $1.32472$ but $W_{29}$'s as a
``parity fit'' at $1.21470$. The recurrent-class series has the exact Padovan
law; the WCC series has only a fit.

\subsection{Select All, Copy, Paste, and the $4^{N/5}$}
\label{sec:copypaste}

$W_{156}/W_{198}$'s largest sector runs
\[
6, 9, 12, 16, 20, 27, 36, 48, 64, 81, 108, 144, 192, 256, 324, 432, 576
\]
over $N=6\ldots22$, with asymptotic growth $4^{N/5}=1.319508^N$. This is
\oeis{A178715} exactly, with no shift.

The entry is Rowell's ``Select All, Copy, Paste'' problem
(\emph{J.\ Integer Seq.} \textbf{18} (2015), 15.10.7): given a budget of $n$
keystrokes, and the ability to type one letter, or to select-all/copy/paste and
thereby multiply the current text by a factor $m$ at a cost of $m+1$ keystrokes,
what is the longest text obtainable? The answer maximises a product of factors
whose \emph{costs} sum to the budget, where factor $m$ costs $m+1$. Optimising
$m^{1/(m+1)}$ over integers gives $m=4$ at $4^{1/5}=1.319508$, beating $m=3$
($3^{1/4}=1.316$) and $m=5$ ($5^{1/6}=1.308$) --- and the entry's own recurrence
$a(n)=4a(n-5)$ for $n\ge16$ is the statement that the optimum is eventually all
fours.

That is where our exponent comes from. The largest Krylov sector of
$W_{156}$ is a product over the gaps between frozen walls, each gap of length
$\ell$ contributing a factor, and the wall itself costing sites; the optimum
therefore has the same shape as the keystroke problem, with the same integer
optimisation and the same answer. R18 derived the Fibonacci recurrence for the
\emph{number} of wall sets from the grammar; A178715 supplies the matching
account of the largest \emph{preimage}.

\subsection{$W_{108}$'s frozen states, factorised by the entry itself}

R19 derived
\[
n_{\rm frozen}(W_{108}, N) \;=\; F_{\lfloor N/2\rfloor+2}\;
F_{\lceil N/2\rceil+2},
\]
a factorisation over the two sublattices. The series is \oeis{A006498}$(N+2)$,
defined by $a(n)=a(n-1)+a(n-3)+a(n-4)$ --- and the entry's own formula section
states the Fibonacci-product identity. So the sublattice factorisation is not
special to this model; it is a known property of that sequence, and our
derivation and the encyclopedia's agree.

$W_{201}$'s frozen count is \oeis{A195971}$(N-1)$, whose stated combinatorial
meaning (a count of $n\times1$ arrays over $0..4$ with a neighbour condition) has
no evident relation to anything here. We record the match and claim nothing
from it.

\section{Two identifications that are new results}

\paragraph{$W_{105}$: central binomial.} The largest sector is
\[
D_{\max}(W_{105}, N) \;=\; \binom{N+1}{\lfloor (N+1)/2\rfloor}
\;=\; \oeis{A001405}(N+1),
\]
verified over $N=6\ldots22$. This is the maximal level set of a balanced
counting statistic, which is exactly what R21 then confirms: $W_{105}$ has a
single conserved range-2 charge, the staggered domain-wall number $S$, and its
level sets \emph{are} the Krylov sectors. The central binomial coefficient is
the largest of those level sets, so the closed form is not a coincidence but the
symmetry showing through.

\paragraph{$W_{150}$: an Euler characteristic.} $D_{\max}(W_{150},N)$ is
\oeis{A214282}$(N+2)$, ``largest Euler characteristic of a downset on an
$n$-dimensional cube'', a unique hit. $W_{150}$'s single conserved charge is the
total domain-wall number $D$ (R21), whose level sets are again exactly its
sectors, so $D_{\max}$ is the largest such level set. The two sequences agree
over the full $17$ terms and diverge from $W_{105}$'s central binomials at
$N=9$ ($210$ against $252$) and $N=13$ ($3003$ against $3432$) --- the two rules
are close but not the same, and the encyclopedia separates them.

R21~\S6 now derives both. $W_{150}$'s sector sizes are $\binom{N+1}{D}$ over
\emph{even} $D$ and $W_{105}$'s are $\binom{N+1}{w}$ over
$w \equiv \lceil N/2\rceil \pmod 2$ --- the same row of Pascal's triangle cut
into its two parity classes. $W_{105}$'s class always contains the centre, hence
the central binomial; $W_{150}$'s misses it exactly at $N \equiv 1 \pmod 4$,
which is where the divergences above sit and where they will keep sitting
($N = 17, 21, \ldots$). The identification came first and the derivation
second, which is the useful direction for an encyclopedia to work in.

\paragraph{And one gap the derivation closed.} With the closed form in hand it
was obvious what $W_{105}$'s \emph{sector count} had to be --- the size of a
parity class, a period-4 staircase --- and it is \oeis{A004525}$(N+2)$, ``one
even followed by three odd''. That row had been missing from the table.

\section{What is not there}

Four series were searched for and returned nothing:

\begin{center}
\small
\begin{tabular}{lll}
\toprule
rules & bc & series \\
\midrule
$W_{28}, W_{70}, W_{157}, W_{199}$ & \rt{pbc} & $n_{\rm rec}$
  $= 1,3,4,3,6,6,8,11,13,18,\ldots$ \\
$W_{29}, W_{71}$ & \rt{pbc} & $n_{\rm rec}$
  $= 1,2,3,2,5,5,7,10,12,17,\ldots$ \\
$W_{28}, W_{70}, W_{157}, W_{199}$ & \rt{pbc} & $D_{\max}$
  $= 1,1,2,1,2,4,2,4,8,4,\ldots$ \\
$W_{29}, W_{71}$ & \rt{pbc} & $D_{\max}$
  $= 2,1,2,1,2,4,2,4,8,4,\ldots$ \\
\bottomrule
\end{tabular}
\end{center}

\noindent The branches are listed separately because the ring splits a group
that \rt{obc0} holds together: at \rt{obc0} all six rules share one class-count
sequence, whereas on the ring $\{28,70,157,199\}$ and $\{29,71\}$ differ at
every $N\ge2$. Their $D_{\max}$ ladders differ only at $N=1$, where the
``ring'' is a single site.

Absence is recorded rather than omitted, because ``we searched and there is
nothing'' is a usable statement. It is also not the same as ``structureless'':
the unmatched \rt{pbc} ladder satisfies $a(N)=2a(N-3)$ exactly from $N=5$, which
\fn{test\_oeis} pins.

\section{Two junk hits, so that they are not rediscovered}

\paragraph{\oeis{A173862}.} The $\rho$-group's \rt{obc0} $D_{\max}$ ladder
$2^{\lfloor (N+1)/3\rfloor}$ has exactly one OEIS home, and it is
``$a(n)=\texttt{A158772}(n-1)/21$''. The terms are right and the match is exact,
but nobody should cite it: the citable object is the closed form, and the
encyclopedia simply has no primary entry for $2^{\lfloor n/3\rfloor}$. We list
it in the table for completeness and use the closed form everywhere else.

\paragraph{The trivial periodics.} $W_{150}$'s frozen count matches
\oeis{A000034} (``period 2: repeat $1,2$''), $W_{105}$'s matches \oeis{A007877}
(``period 4 zigzag $0,1,2,1$''), and $W_{150}/W_{156}$'s sector counts match
\oeis{A008619} (``positive integers repeated''). These are correct and
completely uninformative --- any period-2 or period-4 sequence matches
something. They are kept because the shift is still a check on the closed form:
$W_{105}$'s zero at $N\equiv1\pmod 4$, which R19 had to special-case in
\fn{paper\_figures.\_tier1a}, is visible in A007877's phase.

\section{Reproduce}

\begin{verbatim}
python -m qca_fragmentation.scaling.oeis --refresh --tables
python -m pytest qca_fragmentation/tests/test_oeis.py -q
\end{verbatim}

\noindent \fn{--refresh} is the only step that touches the network; everything
else reads the cached entries in \fn{analytics/}. The identifications, the
partner-agreement checks and the unmatched list are written to
\fn{oeis\_identifications.json} in the same directory.

\end{document}
